%!TeX root=csp.tex
\section{The complexity layer}\label{sec:complexity-layer}

Targets \textup{(T2)} and \textup{(H1)} are stated here and not proved. The
point of stating them is to be precise about what is \emph{not} being claimed.

\begin{definition}[What a cost model would have to supply]\label{def:cost-model}
To state \textup{(T2)} one needs: an encoding of instances of $\CSP(\Gamma)$ as
strings; a machine model with a step count; a predicate ``$f$ is computed in
time bounded by a polynomial in the input length''; and closure of that
predicate under composition. To state \textup{(H1)} one needs, in addition, a
polynomial-time many-one reduction relation $\le_{p}$ and the class
\textup{NP}.
\end{definition}

Mathlib supplies the machine (\texttt{Turing.TM2}), an encoding
(\texttt{Computability.Encoding}), and a time-bounded computation predicate
(\texttt{Turing.TM2ComputableInPolyTime}). It does not supply a complexity
\emph{class}, \textup{P}, \textup{NP}, \textup{NP}-completeness, $\le_{p}$, or
\textsc{Sat}; and the composition lemma
\texttt{TM2ComputableInPolyTime.comp} is an open \texttt{proof\_wanted} in
Mathlib itself. Moreover \texttt{TM2ComputableInPolyTime} applies to total
functions $\alpha\to\beta$, never to languages, so even the existing predicate
would need reshaping.

\begin{remark}\label{rem:complexity-scope}
Building this layer is a substantial project --- comparable in size to the
algebraic core, and with a large formal-verification literature of its own ---
and it contributes none of the mathematics of the dichotomy. Attempting it as
part of this project would be a scoping error. It is exactly the kind of
component that should be built once, in Mathlib, by people interested in
complexity theory, and then consumed.
\end{remark}

\section{Formalization}\label{sec:formalization}

\subsection{What is reused}\label{sec:reuse}

The predecessor development \texttt{ZhukLean} is a complete, \texttt{sorry}-free
Lean~4 formalization of Zhuk's centre theorem, built on
\texttt{Mathlib.ModelTheory}. It is consumed here as a \texttt{lake} dependency,
not copied. It supplies:

\begin{itemize}\tightlist
\item products of structures (\texttt{piStructure}, \texttt{prodStructure},
  coordinatewise realization, reindexing and evaluation homomorphisms) --- the
  one part of the universal-algebra vocabulary Mathlib lacks;
\item absorption in the tuple form (\texttt{Witnesses}, \texttt{Absorbs},
  \texttt{BinAbsorbs}), Taylor terms (\texttt{IsTaylorOn}), and central
  absorption (\texttt{CentrallyAbsorbs});
\item the relational description of absorption by essential relations
  (Barto--Kazda), which is what converts absorption at some arity into
  absorption at arity~$3$;
\item Zhuk's centre theorem itself (\texttt{zhuk\_center}) and the ternary
  collapse (\texttt{exists\_ternary\_witnesses}) --- which is precisely
  Lemma~\ref{lem:central-ternary}, cited without proof in
  \cite{ZhukSimplified}.
\end{itemize}

Three of \cite{ZhukSimplified}'s sixteen black-box imports are therefore already
discharged.

\subsection{What is built here}\label{sec:built}

\begin{center}
\begin{tabular}{@{}llr@{}}
\hline
Module & Contents & Status \\
\hline
\texttt{CSP/Product} & products of structures & complete \\
\texttt{CSP/Congruence} & \texttt{Congruence L M}, order, \texttt{sInf}, quotient & complete \\
\texttt{CSP/Irreducible} & irreducibility, existence and uniqueness of $\cover\sigma$ & complete \\
\texttt{CSP/Bridge} & bridges, collapse, composition, linear/PC & 2 imports open \\
\texttt{CSP/Types} & the six types, multi-types, $\lll$ & complete \\
\texttt{CSP/Instance} & instances, reductions, consistency, cruciality & complete \\
\texttt{CSP/Main} & the main statements & targets \\
\hline
\end{tabular}
\end{center}

Everything below \texttt{CSP/Main} is \texttt{sorry}-free. The two open items in
\texttt{CSP/Bridge} are Lemma~\ref{lem:bridge-comp} and its collapse identity,
both imported from \cite{ZhukJACM}.

\subsection{Design decisions}\label{sec:design}

\begin{description}
\item[Build on \texttt{Mathlib.ModelTheory}, keep the language generic.]
  A bespoke ``finite algebra with one WNU operation'' type would bundle the
  standing hypotheses as fields and avoid threading them, which is a real cost
  in the predecessor development. But it would discard
  \texttt{Substructures.lean} --- 985 lines of exactly the $\Sg$ API needed,
  including \texttt{mem\_closure\_iff\_exists\_term} with the variable type
  equal to the generating set, an interface no hand-rolled clone matches --- and
  it would discard the predecessor's 1600 lines wholesale. Congruences and
  quotients are also \emph{cheaper} in \texttt{ModelTheory}, not dearer.
  Instantiating to $\Vn$ is a thin layer on top.

\item[$\lll$ is \texttt{Relation.ReflTransGen}.]
  Its induction principle is the one the proofs use.

\item[Dotted relations are defined as disjunctions.]
  \texttt{DotLll C B := C = }$\varnothing$\texttt{ }$\vee$\texttt{ Lll C B}, so the case split is forced at every
  use site rather than resolved by the reader.

\item[The bridge criterion is the definition of linearity.]
  Zhuk defines a linear congruence by a per-block isomorphism to
  $\mathbb Z_{p}^{m}$ and derives the bridge criterion. We reverse this: the
  criterion is first-order over finite data and is what every consumer uses; the
  block description becomes a structure theorem. See
  Hazard~\ref{haz:linear-def}.

\item[$\cover\sigma$ is data attached to irreducibility.]
  \texttt{Irreducible.exists\_isCover} produces it; there is no standalone
  function, because there is no cover without irreducibility. A pleasant
  by-product: with the definition taken literally, irreducibility \emph{implies}
  $\sigma\neq A^{2}$, since the empty family intersects to the universe. The
  source never says this.

\item[Instances are lists; scopes are tuples; the variable type is a parameter.]
  See Hazards~\ref{haz:instance-list} and~\ref{haz:renaming}. The last of these
  is the decision that has not yet been cashed out, and it should be made before
  anything about coverings is written.

\item[Nothing is stubbed with \texttt{sorry} at the level of a definition.]
  \texttt{IsConnectedInstance} and \texttt{IsExpandedCovering} are
  \emph{absent} rather than defined as \texttt{sorry}, because a
  \texttt{def~:~Prop~:=~sorry} is a junk constant that downstream proofs can
  silently lean on. Absence forces the design decision; a stub hides it.
\end{description}

\subsection{Order of attack}\label{sec:order}

\begin{enumerate}\tightlist
\item Lemma~\ref{lem:ba-center-implies}. Small, self-contained, used everywhere,
  and adjacent to what is already formalized.
\item The mixed-prime linear algebra of Hazard~\ref{haz:step2-linear-algebra}.
  Independent of everything else; can be done in parallel.
\item Lemma~\ref{lem:bridge-comp}. Twelve source lines of relational algebra;
  closes the two open items in \texttt{CSP/Bridge}.
\item The coproduct design for coverings, then
  Proposition~\ref{prop:covering-basics}.
\item Theorem~\ref{thm:stable-intersection}. The one hard statement in the
  interface, and the sole source of bridges.
\item Lemma~\ref{lem:bridge-from-instance}. Eight hypotheses, ten pages.
\item Theorem~\ref{thm:main-inductive}. Last.
\end{enumerate}

\subsection{Scope}\label{sec:scope}

The calibration point is that the predecessor formalized about $3.5$ printed
pages in about $1670$ lines of Lean --- roughly $480$ lines per page. Excluding
\S\ref{sec:complexity-layer} and \cite{ZhukSimplified}'s \S4, which is an
independent second result and not needed for the dichotomy, the transitive
closure of what must be proved is on the order of $100$--$130$ printed pages,
so $35\,000$--$60\,000$ lines of Lean.

That is a multi-person-year project, and it should be described as one. What is
finished here is its design: the route, the statements, the vocabulary at
formalization granularity, the hazard list, the missing-layer inventory, and a
building development whose foundations are proved and whose summit is stated.
